\section{Propuesta de solución}

Con base en la problemática anteriormente planteada, se propone el desarrollo de un sistema de aplicación web (WebApp) que se integre al ecosistema del Sistema Integral de Producción, con el objetivo de optimizar la gestión de la información y mejorar la eficiencia de los procesos operativos dentro de las áreas de producción y ensamblaje.

La solución propuesta busca atender directamente las limitaciones identificadas, principalmente la dependencia de registros manuales en papel y la existencia de sistemas de información no integrados. Para ello, el sistema permitirá la digitalización de los procesos de captura de datos en planta, reduciendo errores humanos, evitando la pérdida de información y garantizando la disponibilidad de los datos en tiempo real.

Asimismo, el sistema estará diseñado para centralizar la información proveniente de distintas fuentes, permitiendo su integración en una única plataforma. Esto facilitará la consulta, el análisis y la trazabilidad completa de los pedidos a lo largo de todas las etapas del proceso productivo, atendiendo así los requerimientos operativos de visibilidad y control.

Dentro del marco tecnológico, se propone el uso de React.js como librería para el desarrollo de interfaces de usuario dinámicas e intuitivas, facilitando la interacción del personal operativo con el sistema y permitiendo su escalabilidad ante futuros requerimientos. Por otra parte, se plantea el uso de tecnologías del ecosistema .NET, empleando C\# para la lógica de negocio, lo cual permitirá implementar una arquitectura robusta, mantenible y alineada con los recursos tecnológicos disponibles en la organización.

Para la gestión de datos, se propone el uso de PostgreSQL como sistema de bases de datos, debido a su capacidad para manejar grandes volúmenes de información, garantizar la integridad de los datos y ofrecer un alto rendimiento en consultas, lo que resulta fundamental para cumplir con los requerimientos de acceso oportuno y análisis de información.

En conjunto, esta solución permitirá mejorar la administración de la información, agilizar los procesos de seguimiento y control, garantizar la trazabilidad de los pedidos, reducir la dependencia de procesos manuales y fortalecer la toma de decisiones mediante el acceso a información confiable y oportuna.